\section{Опис побудови коду індивідуального завдання. Варіант №7}

\subsection{Об'єктно-орієнтована архітектура та декомпозиція компонентів}

Згідно з індивідуальним Варіантом №7, було розроблено консольну програму для обліку 
фільмів у кінокаталозі. На відміну від процедурного стилю мови C, проєктування архітектури 
виконано із суворим дотриманням принципів успадкування (\textit{Inheritance}), інкапсуляції 
та поліморфізму. \\

Кодову базу розподілено по окремих файлах простору імен \texttt{ad251.gorbunova} за такою 
логічною структурою:
\begin{enumerate}
    \item \textbf{Базовий суперклас \texttt{Media.java}:} 
    Виступає в ролі абстрактного фундаменту предметної області. Він інкапсулює у собі 
    захищене поле \texttt{protected String title} (назва витвору) та базовий віртуальний 
    метод \texttt{printInfo()} для представлення медіа-продукції у консолі.
    \item \textbf{Дочірній клас \texttt{Movie.java}:} 
    Розширює функціональність базового класу за допомогою ключового 
    слова \texttt{extends}. Клас додає специфічні інкапсульовані характеристики кінострічки: 
    приватні поля жанру (\texttt{String genre}), року випуску (\texttt{int year}), 
    тривалості (\texttt{int duration}) та рейтингу (\texttt{double rating}).
    \item \textbf{Сервісний аналізатор масивів \texttt{MovieCatalog.java}:} 
    Компонент, який інкапсулює всередині себе масив посилань на об'єкти базового 
    типу \texttt{Media[]}. Він повністю ізолює логіку обробки складних структур даних 
    (пошук максимуму, ітераційні цикли) від користувача.
    \item \textbf{Головний клас \texttt{Main.java}:} Точка входу в програму (\texttt{main()}), 
    яка координує створення об'єктів та запускає демонстраційні тести.
\end{enumerate}

\subsection{Реалізація механізмів успадкування та поліморфізму в коді}

Під час проектування взаємодії класів було детально опратовано два найважливіші 
інструменти мови Java:
\begin{itemize}
    \item \textbf{Роль ключового слова \texttt{super}:} 
    Конструктор дочірнього класу \texttt{Movie} першим рядком обов'язково викликає 
    конструктор батьківського класу за допомогою інструкції \texttt{super(title)}. Це гарантує 
    правильну послідовну ініціалізацію полів об'єкта в кучі (\textit{Heap}) під час роботи 
    оператора \texttt{new}.
    \item \textbf{Динамічний об'єктний поліморфізм:} 
    Метод \texttt{printCatalog()} класу \texttt{MovieCatalog} виконує ітераційний перебір 
    масиву \texttt{Media[]}. Під час виклику \texttt{item.printInfo()}, завдяки механізму 
    пізнього (динамічного) зв'язування, віртуальна машина Java (JVM) на льоту визначає
     фактичний тип об'єкта в пам'яті й викликає переозначений метод \texttt{@Override} дочірнього 
     класу \texttt{Movie}, а не базового суперкласу.
    \item \textbf{Безпечне приведення типів даних:} 
    Для виконання індивідуального завдання з пошуку фільму з найвищим рейтингом, у 
    класі \texttt{MovieCatalog} застосовано логічний оператор \texttt{instanceof}. Він 
    перевіряє, чи є об'єкт у масиві базових посилань екземпляром класу \texttt{Movie}, після 
    чого виконується безпечне явне приведення типів \texttt{(Movie) item} для доступу до приватного 
    геттера \texttt{getRating()}.
\end{itemize}
